\chapter{The Campaign: A Worked End-to-End Case Study}
\label{ch:casestudy}

Methodology chapters present knowledge in its cleaned-up, logical order.
Real campaigns happen in a different order --- with dead ends, revisions,
and decisions made under uncertainty. This chapter replays ours
chronologically, because the \emph{sequencing} of effort is itself a
transferable skill, and because seeing where the score actually came from
is the best antidote to architecture romanticism.

\section{Phase 0: replicate a known-good solution}

We began not with our own ideas but by \textbf{faithfully replicating the
8th-place solution} (whose author published code and a write-up):
her feature engineering (Chapter~\ref{ch:featurepatterns}), her
multi-branch GRU (Chapter~\ref{ch:rnn}), her walk-forward protocol with
online refit (Chapter~\ref{ch:validation}), down to hyperparameters.

Why replication first is the right opening move:

\begin{itemize}
\item It \textbf{calibrates the harness}: until your pipeline reproduces a
  published number, every measurement it emits is suspect. Ours
  reproduced her tier of performance, which certified the metric
  implementation, the leakage boundaries, and the evaluation loop in one
  stroke.
\item It \textbf{front-loads infrastructure}: the replica forced the
  memmap precompute, the day-batch layout, the online-refit loop --- all
  of which every later experiment reused.
\item It \textbf{establishes the real baseline}: not ``zero'' but ``what
  a gold-tier competitor achieves,'' making every subsequent delta
  meaningful.
\end{itemize}

\section{Phase 1: benchmarks and the first surprises}

A structured benchmark over the model zoo on shared slices produced the
ladder of Chapter~\ref{ch:metric} and three early findings that set
strategy:

\begin{itemize}
\item \textbf{Data beats capacity}: 200 $\to$ 280 training days lifted
  every model; doubling width lifted nothing. Effort re-routed from
  architecture tuning to data logistics.
\item \textbf{Online refit is the lever} ($+0.003$--$0.005$ on RNNs) ---
  and the refit-lr sweep found the plateau-and-cliffs law
  (Chapter~\ref{ch:online}).
\item \textbf{The ensemble wall}: signature and transformer streams,
  despite respectable solo scores, took zero weight in blends. Diversity,
  not solo score, became the admission criterion
  (Chapter~\ref{ch:ensembling}).
\end{itemize}

Standing at $+0.0104$ with a three-stream blend --- roughly 8th-place
tier on our local protocol.

\section{Phase 2: the recalibration}

A check of the actual private leaderboard rewrote the strategy document.
We had been benchmarking against $+0.0112$ --- the 8th place's
\emph{public} score; the private reality: 8th $= 0.0104$, 1st
$= \bm{0.0139}$, with a conspicuous gap between the top three
(0.0132--0.0139) and the pack (4th at 0.0117). Two conclusions:

\begin{itemize}
\item Our local $+0.0104$ was already 8th-tier --- the replication had
  succeeded more fully than assumed.
\item The top three, none of whom published, were separated from the pack
  by a margin ($\sim +0.002$) too large for ensemble polish. They likely
  had \emph{structural} insight the pack lacked. Chasing them with more
  bagging was arithmetic denial; the campaign pivoted to forensics.
\end{itemize}

\begin{keybox}[: read the leaderboard as evidence]
The \emph{shape} of a final leaderboard is data about the problem.
A smooth continuum says grind and variance reduction; a step
discontinuity, like this one's after 3rd place, says a discrete idea
exists that the pack missed. Aim accordingly.
\end{keybox}

\section{Phase 3: the forensic pivot}

In one stretch, the campaign built and ran the instruments of Part~II:
the deconvolution of \resp{8} into the underlying signal
(\S\ref{sec:deconv}, validated by its three sanity checks), the lead--lag
atlas over all 79 features (\S\ref{sec:atlas}), and the two-window
stability gate. Yield, within days:

\begin{itemize}
\item the feature taxonomy (lead / reversal / slow-SMA / market-wide),
  stable at Spearman $+0.97$ across 600 days;
\item the discovery that three features arrive pre-smoothed (SMA-160
  fingerprints) --- an un-smoothing opportunity
  (Chapter~\ref{ch:smoothing});
\item nowcast auxiliary heads (\S\ref{sec:auxtargets}), implemented and
  queued for GPU A/B;
\item and the seeds that upgraded the DRW-style symbolic search from
  $+0.00039$ to $+0.00052$ (Chapter~\ref{ch:selection}).
\end{itemize}

The contrast with Phase 1 is the book's thesis in miniature: a month of
architecture benchmarking moved the blend by $\sim +0.001$; days of
forensics generated the entire remaining research agenda.

\section{Phase 4: importing another competition's playbook}

The DRW crypto winner's write-up was ported wholesale
(Chapter~\ref{ch:selection}): correlation-medoid pruning,
SHAP-consistency, symbolic recombination, with our forensic corrections
(skip pruning on curated pools; seed the search from the atlas). Measured
value on the fixed harness: $+0.0005$ --- and, as a byproduct, the
discovery that the strongest combinations were market-gated reversal
terms, itself a new forensic fact.

The meta-lesson: \textbf{solution write-ups from adjacent competitions
are importable technology}. Reading them is cheap; porting them onto your
own harness --- with your own ablation discipline, because boundary
conditions differ --- is a reliably positive-expected-value activity.

\section{Phase 5: disciplined architecture bets}

Only in the final phase, with forensics in hand, did the campaign return
to architectures --- now as \emph{targeted} bets rather than a zoo tour.
The causal iTransformer and inverted-bridge TimeXer
(Chapter~\ref{ch:attention}) were built because the atlas and the
symbolic search had produced specific evidence for cross-variate gates
and a natural endo/exo split. And rather than buying GPU-weeks on faith,
the synthetic triage world (Chapter~\ref{ch:synthetic}) was built to
decide, at laptop cost, whether either bias is realizable at this SNR
--- with the decision rule fixed in advance.

\section{Where the score actually came from}

The campaign ledger, consolidated. Approximate contributions to the final
local blend of $+0.0104$ (and the queued upside beyond it):

\begin{center}
\small
\begin{tabularx}{\textwidth}{@{}>{\raggedright\arraybackslash}Xr@{}}
\toprule
source & contribution \\
\midrule
replicated baseline (features + ModelR + protocol) & $+0.0076$ \\
online refit, tuned per family & $+0.0016$ \\
recency window (280--400 days vs.\ naive full history) & $+0.0012$ \\
three-stream ensemble over best-per-family & $+0.0012$ \\
engineered blocks (rsig, clusters, symbolic combos)$^{\ast}$ &
$+0.0005$--$0.0009$ \\
architecture changes beyond the replica & $\approx 0$ \\
\midrule
queued: nowcast heads, atlas-seeded features at scale, AE factors,
triaged attention & unknown \\
\bottomrule
\end{tabularx}
\end{center}
{\footnotesize $^{\ast}$measured on the XGB stream; blend-level
contribution partially overlapping with ensemble line.}

Read the zero in that table again: \emph{every} architecture beyond the
replicated one --- transformers, signatures, width doublings ---
contributed nothing to the blend that mattered. The score came from
process (adaptation, recency, averaging), features, and calibration.
That is not an argument against architecture research; it is an argument
for sequencing it \textbf{last}, behind forensics, on specific evidence,
through cheap triage.

\begin{drillbox}
Write the phase plan for your next competition before touching data, with
exit criteria per phase: (0) replicate something published --- exit when
your harness reproduces its number; (1) benchmark families on a fixed
harness --- exit at the ensemble wall; (2) recalibrate against the true
target --- exit with a gap analysis; (3) forensics --- exit with a
stability-gated taxonomy; (4) import adjacent playbooks; (5) targeted
architecture bets through synthetic triage. Then notice how much of the
usual first month --- architecture shopping --- this plan defers, and
where it spends the time instead.
\end{drillbox}
